\documentclass{article}
\usepackage{amsmath}
\usepackage{amssymb}
\usepackage{amsthm}
\usepackage{mathtools}
\usepackage{geometry}
\usepackage{hyperref}

\newcommand{\rk}{\operatorname{rank}}
\newcommand{\GL}{\operatorname{GL}}
\newcommand{\R}{\mathbb{R}}
\newcommand{\Mat}{\mathrm{Mat}}

\theoremstyle{definition}
\newtheorem{theorem}{Theorem}[section]
\newtheorem{proposition}[theorem]{Proposition}
\newtheorem{lemma}[theorem]{Lemma}
\newtheorem{corollary}[theorem]{Corollary}
\newtheorem{definition}[theorem]{Definition}
\newtheorem{remark}[theorem]{Remark}

\title{Symmetry Groups and Orbit Spaces}
\date{}

\begin{document}
\maketitle

\section{Groups and Orbit Spaces}

To formulate the first conjecture correctly, we must distinguish two different symmetry groups acting on the parameter space.

\subsubsection*{The full base-change group}

Let
\[
G_d:=\prod_{\ell=0}^L \GL(d_\ell,\mathbb R).
\]
It acts on $\Omega$ by
\[
(g\cdot W)_\ell := g_\ell W_\ell g_{\ell-1}^{-1}.
\]
Under this action,
\[
\mu(g\cdot W)=g_L\mu(W)g_0^{-1}.
\]

\subsubsection*{The hidden-layer gauge group}

Let
\[
G_{\mathrm{hid}}:=\prod_{\ell=1}^{L-1}\GL(d_\ell,\mathbb R),
\]
viewed as the subgroup of $G_d$ with $g_0=I_{d_0}$ and $g_L=I_{d_L}$.

Its action on $\Omega$ is
\[
g\cdot W
=
\bigl(
W_Lg_{L-1}^{-1},
\,
(g_\ell W_\ell g_{\ell-1}^{-1})_{2\le \ell\le L-1},
\,
g_1W_1
\bigr).
\]
For this action,
\[
\mu(g\cdot W)=\mu(W).
\]

\begin{remark}[Why the hidden gauge group is too small for the main conjecture]
Because $G_{\mathrm{hid}}$ fixes the end-to-end map $\mu(W)$ exactly, two points can lie in the same $G_{\mathrm{hid}}$-orbit only if they already have the same product.

Therefore $G_{\mathrm{hid}}$ is the natural symmetry group of the critical locus itself, but it is too small to identify a general point of $\Gamma_{d,S}^r$ with a point of $\mathcal C_{d,S}^r$ unless the equation
\[
\mu(W)=P_SW^\ast
\]
is already built into the definition of the source space.
\end{remark}

\subsubsection*{The stabilizer of the closing map}

Fix
\[
W_0=W_S:=\Sigma_{XY}P_Q\in \Mat_{d_0\times d_L}(\mathbb R).
\]
Define the stabilizer subgroup
\[
G_d^{W_0}
:=
\{g\in G_d : g_0W_0 g_L^{-1}=W_0\}.
\]

\begin{proposition}
The weight-space cyclic locus $\Gamma_{d,S}$ is stable under the action of $G_d^{W_0}$.
\end{proposition}

\begin{proof}
Let
\[
E_\ell(W):=W_{<\ell}W_0W_{>\ell}.
\]
For $g\in G_d^{W_0}$ one has
\[
\qquad
(g\cdot W)_{<\ell}=g_{\ell-1}W_{<\ell}g_0^{-1},
\qquad
(g\cdot W)_{>\ell}=g_LW_{>\ell}g_\ell^{-1},
\]
hence
\[
E_\ell(g\cdot W)
=
g_{\ell-1}W_{<\ell}(g_0^{-1}W_0g_L)W_{>\ell}g_\ell^{-1}.
\]
Since $g_0W_0g_L^{-1}=W_0$, equivalently
\[
g_0^{-1}W_0g_L=W_0,
\]
we obtain
\[
E_\ell(g\cdot W)=g_{\ell-1}E_\ell(W)g_\ell^{-1}.
\]
Therefore $E_\ell(W)=0$ for all $\ell$ if and only if $E_\ell(g\cdot W)=0$ for all $\ell$.
\end{proof}

\begin{remark}
The same computation shows that $\Gamma_{d,S}^r$ is stable under $G_d^{W_0}$ because the rank of the end-to-end map is preserved under left-right multiplication by invertible endpoint matrices.
\end{remark}

\subsubsection*{Orbit spaces relevant to the setup}

The first orbit spaces to consider are
\[
\Gamma_{d,S}^r / G_d^{W_0},
\qquad
\mathcal C_{d,S}^r / G_d^{W_0}.
\]
These are the natural candidates for the first conjecture because $G_d^{W_0}$ can move the endpoint map while preserving the closing matrix $W_0$ that defines the cyclic equations.

\begin{remark}[A practical reformulation of the first conjecture]
The desired orbit correspondence can be phrased as follows:
\begin{quote}
Every orbit of $G_d^{W_0}$ in $\Gamma_{d,S}^r$ contains at least one point whose end-to-end map is $P_SW^\ast$, and two such points lie in the same $G_d^{W_0}$-orbit inside $\Gamma_{d,S}^r$ if and only if they lie in the same $G_d^{W_0}$-orbit inside $\mathcal C_{d,S}^r$.
\end{quote}
If this holds, then the inclusion
\[
\mathcal C_{d,S}^r \hookrightarrow \Gamma_{d,S}^r
\]
induces a bijection on orbit spaces.
\end{remark}

\begin{remark}[Open point for later verification]
The exact choice of acting group may still need refinement. One could imagine replacing $G_d^{W_0}$ by a subgroup that fixes not only $W_0$ but also the singular subspace decomposition attached to $S$. For the moment, however, $G_d^{W_0}$ is the cleanest candidate for the first direct weight-space conjecture.
\end{remark}

\end{document}
